% !TeX root = ../../Book3.tex
% 纲要标题：### 22.4 Gorenstein 环的对偶例子
% 纲要位置：计划纲要.md，第 2692 行。

\section{Gorenstein 环的对偶例子}
\label{b3:sec:2204}

零维 Gorenstein 性可以通过基座及乘法配对检验，因而适合具体计算。正维及非 Cohen–Macaulay 情形还要保留对偶的次数；一个典范模未必足以记录全部数据。


% 纲要内容（仅作写作计划，尚非教材正文）：
% * Artin Gorenstein 环的基座与乘法配对
% * 超曲面及完全交的环境环 Ext 计算
% * 环境环上的对偶复形：有限长度和各项有限秩的自由消解、对偶微分及移位；区分各项的 $S$ 模结构与上同调的 $R$ 模结构
% * 非 CM 例子的两个 Ext 次数与局部上同调的 Matlis 对应
% ---

\subsection{Artin Gorenstein 环}
\label{b3:subsec:2204:01}

零维时，Gorenstein 性可以通过一个很小的线性空间辨认。基座包含那些被极大理想全部零化的元素；它们是乘法再也不能向下推进的部分。

\begin{proposition}\label{b3:prop:artinian-gorenstein-socle}
设 \((R,\mathfrak m,k)\) 为非零交换 Artin 局部环，并记 \(\operatorname{Soc}(R)=(0:_R\mathfrak m)\)。则下列条件等价：\(R\) 为 Gorenstein 环；\(R\) 为内射 \(R\) 模；\(\dim_k\operatorname{Soc}(R)=1\)。
\end{proposition}
\begin{proof}
零维时，\cref{b3:thm:gorenstein-bass-criterion}给出内射维数为零，并使
\(\operatorname{Hom}_R(k,R)=\operatorname{Soc}(R)\) 为一维。下面也直接说明基座条件如何产生自内射性。Artin 局部环完备，其有限长度模 \(R\) 的 Matlis 对偶 \(D_R(R)=E_R(k)\) 与 \(R\) 有相同长度；这是\cref{b3:thm:matlis-duality}在组成列上的应用，因为 \(D_R(k)\cong k\) 且对偶正合。

若基座一维，把 \(\operatorname{Soc}(R)\cong k\) 嵌入 \(E_R(k)\)，利用内射性将它延拓为 \(u:R\to E_R(k)\)。任一非零理想 \(J\) 都与基座非零相交：因 \(\mathfrak m\) 幂零，选最大 \(r\) 使 \(\mathfrak m^rJ\ne0\)，则该模被 \(\mathfrak m\) 零化。因此 \(\ker u\) 若非零，必与基座相交，与 \(u\) 在基座上单射矛盾。故 \(u\) 单射，两端长度相同使它为同构，\(R\) 遂内射。
\end{proof}

若 \(R=A\) 还是剩余域为 \(k\) 的有限维交换局部 \(k\) 代数，上述结论可直接写成乘法配对。选 \(k\) 线性泛函 \(\lambda:A\to k\)，使它在一维基座上非零。则
\[
A\longrightarrow\operatorname{Hom}_k(A,k),\qquad
a\longmapsto\bigl(b\longmapsto\lambda(ab)\bigr)
\]
单射：若 \(a\ne0\)，理想 \(Aa\) 包含非零基座元素，所以存在 \(b\) 使 \(\lambda(ab)\ne0\)。两端维数相同，故这是同构。这里的目标模取作用 \((c\phi)(b)=\phi(cb)\)；它内射，因为
\[
\operatorname{Hom}_A(N,\operatorname{Hom}_k(A,k))\cong\operatorname{Hom}_k(N,k).
\]
同构把 \(f\) 送到 \(n\mapsto f(n)(1)\)，逆映射把 \(u\) 送到 \(n\mapsto[a\mapsto u(an)]\)。向量空间对偶正合，从而此 Hom 函子正合。

\ProseParagraphBreak
例如 \(A=k[x,y]/(x^a,y^b)\)、\(a,b\ge1\)，基座由 \(x^{a-1}y^{b-1}\) 生成。取 \(\lambda\) 为此单项式的系数，基单项式 \(x^iy^j\) 与 \(x^{a-1-i}y^{b-1-j}\) 配对为一，其余达到不同次数的乘积对该系数贡献为零。这个有限维配对把“类型一”的条件变成能够逐项检验的矩阵。同样长度很短的 \(k[x,y]/(x,y)^2\) 却有二维基座 \(kx+ky\)，所以不存在这样的非退化乘法配对。

\ProseParagraphBreak
对后一环 \(A=k[x,y]/(x,y)^2\)，典范模仍然存在，但不再同构于环本身。令 \(W=\operatorname{Hom}_k(A,k)\)，以 \(\phi_0,\phi_x,\phi_y\) 表示基 \(1,x,y\) 的对偶基。按照 \((a\phi)(b)=\phi(ab)\) 的作用，直接计算得
\[
\begin{aligned}
x\phi_x&=\phi_0,& y\phi_y&=\phi_0,\\
x\phi_y&=0,& y\phi_x&=0,\\
x\phi_0&=0,& y\phi_0&=0.
\end{aligned}
\]
于是 \(\operatorname{Soc}_A(W)=k\phi_0\)。任意非零子模中，若一个元素的 \(\phi_x\) 或 \(\phi_y\) 系数非零，乘 \(x\) 或 \(y\) 就得到非零的 \(\phi_0\)；否则它本身已在 \(k\phi_0\) 中。因此 \(k\phi_0\subseteq W\) 本质。上面的 Hom 计算又说明 \(W\) 内射，故 \(W\cong E_A(k)=D_A(A)\)。零维时 \(H^0_{\mathfrak m}(A)=A\)，Matlis 双对偶遂使 \(W\) 成为典范模。可是 \(\mathfrak mW=k\phi_0\)，所以 \(W/\mathfrak mW\) 为二维，\(W\) 至少需要两个生成元；\(A\) 作为自身模只需一个。这个作用表具体说明了 \(\omega_A\not\cong A\)，也区分了典范模的存在与 Gorenstein 性。

\subsection{超曲面与完全交的对偶例子}
\label{b3:subsec:2204:02}

设 \(k\) 为域，取 \(S=k[\![x,y]\!]\)、\(R=S/(xy)\)。环境环中的自由消解是
\[
0\longrightarrow S\xrightarrow{xy}S\longrightarrow R\longrightarrow0.
\]
对偶后仍为乘 \(xy\)，所以 \(\omega_R=\operatorname{Ext}_S^1(R,S)\cong R\)。这是一维 Gorenstein 环，但在原点有两条分支，因而不是整环，也不是正则局部环。Gorenstein 条件描述对偶性质，不能替代不可约性或光滑性。

\ProseParagraphBreak
该例的局部上同调也可以算出。因为在 \(R\) 中 \(xy=0\)，有 \(R_{xy}=0\)，所以两元 Čech 复形给出
\[
H_{\mathfrak m}^1(R)=\bigl(R_x\oplus R_y\bigr)/R
=\bigl(k((x))\oplus k((y))\bigr)/R.
\]
此处 \(R\) 在两支中的像是所有常数项相等的幂级数对：任一元素唯一写成
\(c+x\cdot a(x)+y\cdot b(y)\)。故上式中的 \(R\to R_x\oplus R_y\) 为单射，\(H_{\mathfrak m}^0(R)=0\)。局部对偶进一步给出
\(D_R(H_{\mathfrak m}^1(R))\cong R\)，并由 Matlis 双对偶得到
\(H_{\mathfrak m}^1(R)\cong E_R(k)\)。最后这个同构反映的是内射包络结构，不应被误写成 \(H_{\mathfrak m}^1(R)\cong R\)。

\ProseParagraphBreak
再设 \(k\) 为域，取 \(S=k[\![x,y,z]\!]\)、\(R=S/(x^2,y^3)\)。以 \((x^2,y^3)\) 为序，消解为
\[
0\longrightarrow S\xrightarrow{\binom{-y^3}{x^2}}S^2
\xrightarrow{(x^2\ \ y^3)}S\longrightarrow R\longrightarrow0.
\]
取对偶后，零次核为零；若 \(-y^3a+x^2b=0\)，则互素性给出
\((a,b)=(x^2c,y^3c)\)，故一次上同调为零；二次余核为 \(R\)。所以 \(\omega_R\cong R\)，且 \(\dim R=1\)。参数 \(z\) 为非零因子，商去 \(z\) 后得到前面有非退化配对的 Artin 环 \(k[x,y]/(x^2,y^3)\)。

\ProseParagraphBreak
这些对偶计算全部在环境环 \(S\) 上进行。商环 \(R\) 上的剩余域自由消解通常无限长；例如 \(k[\varepsilon]/(\varepsilon^2)\) 上由乘 \(\varepsilon\) 构成的无限消解已在\secref{b3:sec:1903}计算，而环本身却内射。必须在记号中保留 \(\operatorname{Ext}_S\)、\(\operatorname{Ext}_R\) 的下标，才能区分这两项事实。

\subsection{环境环上的对偶复形}
\label{b3:subsec:2204:03}

对完备正则局部环 \(S\) 的一个 CM 商 \(R=S/I\)，自由消解的对偶只有一个非零上同调模，因而典范模能够集中表达对偶信息。若 \(R\) 不是 CM 环，不同次数上的对偶信息可能同时存在；这时保留整个对偶复形，比只取最高次局部上同调的对偶更合适。

\begin{definition}\label{b3:def:regular-cover-dual-complex}
设 \((S,\mathfrak n,k)\) 为 \(n\) 维交换完备 Noether 正则局部环，\(I\subsetneq S\) 为理想，\(R=S/I\)。取有限长度、各项均为有限秩自由 \(S\) 模的消解 \(P_\bullet\rightarrow R\)，微分记为 \(\mathrm d_j:P_j\rightarrow P_{j-1}\)。令
\(C^j=\operatorname{Hom}_S(P_j,S)\)，微分为
\(\delta^j(\varphi)=\varphi\circ\mathrm d_{j+1}\)。把
\[
\operatorname{Hom}_S(P_\bullet,S)[n]
\]
看作上链复形，称它为给定满射表示 \(S\twoheadrightarrow R\) 下的一个归一化对偶复形模型。移位约定为
\(C[n]^j=C^{j+n}\)、\(\mathrm{d}_{C[n]}=(-1)^n\delta\)，所以对每个整数 \(i\)，该模型的上同调为
\[
H^{-i}(C[n])=\operatorname{Ext}_S^{n-i}(R,S).
\]
未出现的 \(P_j\) 及负次 Ext 均记为零。
\end{definition}

上述有限消解可以取得：\cref{b3:thm:regular-local-global-dimension}给出正则局部环 \(S\) 的整体维数为 \(n\)，而有限生成模的极小自由消解在每一项只有有限个生成元。若 \(P_\bullet\) 与 \(Q_\bullet\) 都是这样的消解，恒等映射 \(R\rightarrow R\) 可提升为两个方向的比较映射；其复合与各自的恒等映射链同伦。取 \(S\) 对偶并移位后，仍得到互为逆的上链同伦等价。因此这里的消解选择不改变模型的链同伦类型，也不改变各次上同调。

\ProseParagraphBreak
该模型的各项是 \(S\) 模，未必是 \(R\) 模，因为 \(I\) 未必零化这些自由项。不过，对任意 \(a\in I\)，消解上的乘 \(a\) 映射诱导 \(R\) 上的零映射，比较映射的同伦唯一性使它链同伦于零。对偶后乘 \(a\) 同样同伦于零，故 \(I\) 零化每个上同调模。这才给出上述 Ext 模的 \(R\) 模结构。

\ProseParagraphBreak
若 \(R\) 为 \(d\) 维 CM 环，\cref{b3:thm:regular-local-duality}及局部上同调的消失说明只有
\(\operatorname{Ext}_S^{n-d}(R,S)=\omega_R\) 非零。因此归一化后的模型在次数 \(-d\) 的上同调为 \(\omega_R\)，其余次数均为零；这个次数与局部对偶公式中的 \(d-i\) 相对应。

\ProseParagraphBreak
最后设 \(k\) 为域，取 \(S=k[\![x,y]\!]\)、\(R=S/(x^2,xy)\)，并记 \(\mathfrak m=(x,y)R\)。若 \(x^2a+xyb=0\)，在整环 \(S\) 中消去 \(x\) 得 \(xa+yb=0\)；由于 \(x,y\) 互素，必有 \((a,b)=c(-y,x)\)。故有自由消解
\[
0\longrightarrow S\xrightarrow{\binom{-y}{x}}S^2
\xrightarrow{(x^2\ \ xy)}S\longrightarrow R\longrightarrow0.
\]
取 \(S\) 对偶后，两个微分分别是
\[
S\xrightarrow{\binom{x^2}{xy}}S^2
\xrightarrow{(-y\ \ x)}S.
\]
若 \(-ya+xb=0\)，同样的整除论证给出 \((a,b)=c(x,y)\)。所以一次核由 \((x,y)\) 生成，而前一像由 \(x(x,y)\) 生成；二次余核则是 \(S/(x,y)\)。因此
\[
\operatorname{Ext}_S^1(R,S)\cong S/(x),\qquad
\operatorname{Ext}_S^2(R,S)\cong S/(x,y)=k.
\]
这里 \(S\) 的维数为二，归一化对偶模型在次数 \(-1\) 和零都有非零上同调。\subsecref{b3:subsec:2201:03}给出的两个局部上同调模正好与之对应。具体地，\(R\) 为完备 Noether 局部环，取
\(D_R(-)=\operatorname{Hom}_R(-,E_R(k))\)。由\cref{b3:thm:regular-local-duality}及商环上的对偶识别，得到\cref{b3:tab:non-cm-dual-degrees}中的两项同构；Ext 的下标仍是环境环 \(S\)。

\begin{table}[htbp]
\centering
\caption{非 Cohen--Macaulay 环的两个对偶次数}
\label{b3:tab:non-cm-dual-degrees}
\begin{tabularx}{\linewidth}{@{}cXX@{}}
\toprule
次数 \(i\) & 局部上同调 \(H^i_{\mathfrak m}(R)\) & Matlis 对偶 \(D_R(H^i_{\mathfrak m}(R))\) \\
\midrule
\(1\) & \(k((y))/k[\![y]\!]\) & \(\operatorname{Ext}_S^1(R,S)\cong S/(x)\) \\
\(0\) & \(kx\cong k\) & \(\operatorname{Ext}_S^2(R,S)\cong k\) \\
\bottomrule
\end{tabularx}
\end{table}

典范模 \(\omega_R=\operatorname{Ext}_S^1(R,S)\) 记录了一次局部上同调的对偶，却没有记录由 \(kx\) 产生的另一个非零次数。这个例子说明，非 CM 环的局部对偶信息通常需要多个模及其次数，不能仅由一个典范模替代。
